\documentclass[11pt]{article}
\usepackage[margin=25mm]{geometry}
\usepackage{fontspec}
\setmainfont{Latin Modern Roman}
\usepackage{amsmath,amssymb,mathtools}
\usepackage{unicode-math}
\setmathfont{Latin Modern Math}
\usepackage{xurl,hyperref}
\hypersetup{colorlinks=true,urlcolor=blue}
\setcounter{secnumdepth}{0}
\setlength{\parindent}{0pt}
\setlength{\parskip}{5pt}
\setlength{\emergencystretch}{3em}
\newcommand{\sref}[2]{\hyperlink{source:#1:#2}{[S#2]}}
\newcommand{\cataloguescope}[1]{\paragraph{Recorded target.}#1}
\newcommand{\R}{\mathbb R}
\newcommand{\C}{\mathbb C}
\newcommand{\N}{\mathbb N}
\newcommand{\Z}{\mathbb Z}
\newcommand{\Q}{\mathbb Q}
\newcommand{\F}{\mathbb F}
\newcommand{\eps}{\varepsilon}
\begin{document}
\section{Lower Bounds for Hamiltonian Mass}\label{statement}

\subsection{Definitions and mathematical statement}

\paragraph{Shared notation.}$\mathbb N=\{1,2,\ldots\}$; the standard number fields have their usual meanings. All additional conventions are stated locally.


An initial data set consists of a spatial metric, second fundamental form and matter fields satisfying the gravitational constraints. A dominant-energy condition bounds matter momentum by energy density. Hamiltonian mass is the boundary charge generated by an asymptotic time symmetry in the source's asymptotic class.

\paragraph{Catalogue wording retained for traceability.}
For every globally well-behaved initial data set whose matter satisfies the dominant-energy-type positivity assumptions, prove that its Hamiltonian mass is bounded below under asymptotic symmetries.

\paragraph{Scope and source qualification.}Globally well-behaved and the allowed asymptotic symmetries require the report's precise hypotheses; asymptotically flat and asymptotically anti-de Sitter charges should not be interchanged.
\subsection{Short English statement}
Prove lower bounds on the Hamiltonian mass of globally well-behaved initial data satisfying the stated positivity assumptions.
\subsection{Research attempt}
\paragraph{Model, date and investigated scope.}
OpenAI GPT-6 Astra (configured model \texttt{gpt-6-astra}), ultra reasoning; generated 11 October 2026. Derived radial mass monotonicity for a time-symmetric spherical asymptotically flat test class; no general Hamiltonian positivity theorem.

\paragraph{Sources and literature check.}
Read the full supplied definition, catalogue wording, source qualifications and listed references for this question. No fresh status or novelty search was performed. The argument below proves its auxiliary statements directly and does not import an unverified external theorem.

\paragraph{A radial time-symmetric test.} On a three-dimensional warped metric \(g=ds^2+r(s)^2d\Omega^2\), define the radial mass \(m(s)=r(1-r'^2)/2\). The radial and tangential sectional curvatures are \(-r''/r\) and \((1-r'^2)/r^2\), obtained from the warped metric connection; summing the three sectional curvatures twice gives
\[
 R(g)=\frac{2(1-r'^2-2rr'')}{r^2},\qquad
 m'(s)=\frac{r^2r'}4R(g).
\]
The second identity follows by differentiating the definition of \(m\), so it does not require a mass-positivity theorem.
\paragraph{What the monotonicity proves.} If \(r'>0\), \(R(g)\geq0\), and the center is regular with \(r(0)=0\), \(r'(0)=1\), then \(m(0)=0\) and \(m(s)\geq0\). If a finite asymptotically flat mass limit exists in this radial class, its value is nonnegative. Starting instead at a minimal sphere \(r'=0\) gives \(m=r/2\) there and the same outward monotonicity on a region with \(r'>0\).
\paragraph{Scope limitation.} Time symmetry, spherical symmetry, a monotone area radius and asymptotically flat normalization are strong restrictions. With \(r'<0\) the monotonicity changes sign; with negative cosmological constant the radial charge needs a background term. General second fundamental form and momentum constraints are absent. Thus no lower bound for every globally well-behaved source initial data set or every allowed Hamiltonian asymptotic generator is established. The first gap is a coercive identity for the actual nonsymmetric constraint and asymptotic class, with all interior and boundary terms controlled.

\paragraph{Verification and adversarial audit.}
Checked scalar-curvature coefficients, regular-center and minimal-sphere values, r prime sign, and noninterchangeability of flat and AdS charges. Model self-audit only; no independent human mathematical verification. Applied contribution prompt A before drafting and B afterwards. Checked scalar-curvature coefficients, regular-center and minimal-sphere values, r prime sign, and noninterchangeability of flat and AdS charges. First remaining gap: Construct a coercive identity in the source asymptotic class without spherical/time-symmetric restrictions and control global boundary terms. The inherited source status is preserved; original-target completion remains zero.

Coordinator checked the warped scalar curvature and radial mass derivative under the monotone-radius hypothesis. This additional check was model-assisted, not independent human verification.

\paragraph{Tools, process and contributions.}
Python 3 and local file tools used to read definitions and serialize manuscripts. Calculations are symbolic reasoning; no numerical experiment or theorem prover. Contribution prompt A was followed by prompt B. GPT-6 Astra with ultra reasoning drafted the research and performed its stated self-audit; any additional coordinator contribution is identified in the verification and audit. A GPT-6-based Codex coordinator assembled the manuscript and checked repository integration; its exact variant and reasoning configuration were not exposed. No human mathematical contribution, independent human proof check or formal proof certification is supplied. Complete derivations and the consulted source references are included here; private reasoning transcripts are not published.

\paragraph{Final status and first remaining gap.}
\textbf{UNRESOLVED}. Construct a coercive identity in the source asymptotic class without spherical/time-symmetric restrictions and control global boundary terms.

\section{Completion Estimate}
\noindent\textbf{COMPLETION: 0\%}
This is an elementary auxiliary calculation or an explicit test of a restricted strategy. No new portion of the full source target is established.

\subsection{Sources}
\begin{itemize}
\item[\hypertarget{source:statement:1}{S1}] ULAM.AI and the UnsolvedMath Contributors, UnsolvedMath, catalogue entry. Catalogue curation under CC BY 4.0; definitions and English exposition added locally.
\url{https://huggingface.co/datasets/ulamai/UnsolvedMath}.
\item[\hypertarget{source:statement:2}{S2}] Analysis, Geometry and Topology of Positive Scalar Curvature Metrics (2024), Oberwolfach Reports, DOI 10.4171/owr/2024/9, pp. 483--566. Original report and mathematical context. Incorporated report material remains under the source CC BY-SA 4.0 terms unless otherwise noted.
\url{https://doi.org/10.4171/owr/2024/9}.
\end{itemize}
\label{end:statement}
\end{document}
